\section{Related Work}
\label{app:related-work}

This appendix expands the positioning sketched in Section~\ref{sec:introduction}. The organizing principle is the layer used for comparison. Much of the neural-network verification literature is naturally solver-centered, geometry-centered, or proof-format-centered. Our paper instead studies a proof-centered semantic cover layer, defined only after fixing a verifier-facing language: one counts verifier-visible local certified cells together with an accepted global cover, and then compares proof languages by the overhead with which they realize that layer. The point of this appendix is therefore not simply to list neighboring topics, but to make explicit where the present work sits relative to them.

\subsection{Field Positioning and Solver-Centered Verification}

Broad surveys and taxonomies usually organize neural-network verification by solver family or algorithmic workflow. Representative references include the survey of Liu et al.~\cite{Liu2021Survey}, the unified CPWL-verification view of Bunel et al.~\cite{Bunel2018UnifiedView}, and the more recent programming-languages positioning of Cordeiro et al.~\cite{Cordeiro2025PLChallenge}. These works are important context for the present paper because they clarify the dominant categories in the field: SMT, mixed-integer optimization, branch-and-bound, relaxation, propagation, and reachability. Our paper agrees that these distinctions matter algorithmically, but asks a different question: once a proof-producing pipeline is in place, what is the semantic size of the proof object that the trusted verifier must actually consume?

That question is deliberately orthogonal to the main complete-verification lines. Classical SMT-style and phase-search methods for CPWL networks include Reluplex and the work of Ehlers~\cite{Katz2017Reluplex,Ehlers2017PiecewiseLinear}. Mixed-integer formulations provide another exact route, represented by Tjeng et al.~\cite{Tjeng2019MIP}. Branch-and-bound and relaxation-guided complete verification have since become a major line of development, with representative systems and analyses in \cite{Bunel2020BranchBound,Xu2021FastComplete,Ferrari2022MultiNeuronBaB,Wang2021BetaCROWN,Zhang2022GCPCROWN,Shi2025GenBaB,Chiu2025SDPCROWN}. These works typically measure difficulty through runtime, branching quality, relaxation strength, node counts, or encoding size. By contrast, our VAC formalism abstracts away from the internal search trace and counts verifier-accepted local cells instead.

A similar contrast applies to propagation, abstract-interpretation, and reachability-based verification. Symbolic interval analysis, abstract interpretation, abstract domains, linear bound propagation, and reachable-set tools all produce certified information, but the natural performance proxies in those literatures remain bound tightness, propagation cost, refinement behavior, or set representation size rather than the semantic size of the final proof object~\cite{Wang2018ReluVal,Gehr2018AI2,Singh2019AbstractDomain,Zhang2018CROWN,Lopez2023NNV2}. Framework papers such as Marabou and DNNV are also highly relevant because they make tool-chain decomposition explicit: they separate network representation, property language, solver backend, and proof or checking infrastructure~\cite{Katz2019Marabou,Shriver2021DNNV}. Our paper can be read as pushing that decomposition one step further by asking what the trusted backend ought to count.

In this sense, the present work is not a competing solver architecture. Sections~\ref{sec:verifier-model} through \ref{sec:vac-bounds} aim to define a quantity that can sit above many solver families. Different verifiers may expose different candidate cells, different proof languages, or different local payload schemas, but the semantic question remains the same: how many verifier-visible local units are needed before the global statement can be accepted?

\subsection{Proof Production, Proof Checking, and Certificate Reuse}

The closest line in spirit is proof-producing neural verification. Isac et al.~\cite{Isac2022ProofProduction} show how a verifier can emit independently checkable proof artifacts rather than only final verdicts. Fischer et al.~\cite{Fischer2022SharedCertificates} study how certificates can be shared across related verification instances. More recent work studies certified proof checking inside theorem-proving environments~\cite{Desmartin2025ImandraChecker}, cross-tool proof formats and proof checkers~\cite{Duong2025APTP}, and compositional assume-guarantee reasoning for neural verification~\cite{Duong2025CoVeNN}.

Our contribution is complementary to this literature. Those works ask how to produce, represent, exchange, certify, or reuse proofs. We ask what the semantic size of a proof should be once one looks only at verifier-accepted local units. This is why Section~\ref{sec:vac-proof-languages} separates VAC from proof syntax: within a fixed verifier-facing language, VAC is the semantic cover quantity; SCC is the size of a particular cell-complete proof-language realization; and DL is its encoding length. Put differently, proof production and proof checking make verifier-visible proof objects possible; our paper studies the complexity of those objects after the verifier interface has been fixed.

This distinction also explains why the current paper is language-relative without being syntax-first. APTP-style or theorem-prover-style proof formats are important realizations, but the present paper does not argue that one such format is fundamental. Instead, after fixing the verifier-facing language, it treats the accepted local cell as the semantic proof unit and compares proof languages by how efficiently they realize covers of such units. That is exactly the role played by the cell-complete realization theorems in Section~\ref{sec:vac-proof-languages}.

\subsection{CPWL Geometry, Linear Regions, and Exact Abstraction}

A second nearby literature studies the geometry and expressive power of piecewise-linear networks through linear regions and related decompositions~\cite{Montufar2014LinearRegions,Raghu2017ExpressivePower,Serra2018Counting}. These papers are indispensable background because they explain why CPWL models admit finite hyperplane structure at all. However, they do not by themselves define proof complexity. A verifier-admissible exact cell may be coarser than a geometric region, finer than a geometric region, or overlapping in ways that arrangement cells do not capture. On the threshold side, a proof may not need to reconstruct the full arrangement at all. This is why Section~\ref{subsec:arrangement-degeneration} treats arrangement count only as a restricted exact-side degeneration under connectedness and language-realizability assumptions, rather than as the semantic cover quantity of the theory.

More closely related are works on abstraction with locally exact structure or exact refinement. Manino et al.~\cite{Manino2023GlobalAbstractions} study global abstractions with locally exact reconstruction, and Liu et al.~\cite{Liu2024AbstractionRefinement} develop abstraction-refinement schemes aimed at scalable exact verification. These works are adjacent to our guard-certified exactification results because they also take seriously the idea that local exact structure can be the right bridge between continuous semantics and verifiable reasoning. The difference is in the role played by that local structure. In our paper, exactification is not introduced as another heuristic verification strategy; it is a structural theorem for the CPWL expression graphs under study. Section~\ref{sec:guard-exactification} uses it to prove the existence of finite accepted exact certified covers, and Section~\ref{sec:vac-bounds} then turns that existence result into exact-side upper bounds for a verifier-aligned proof-complexity quantity.

The same distinction matters for candidate generation versus cover complexity. Abstraction and refinement papers often have to solve both problems simultaneously: how to generate good local candidates, and how to reason over them. Our dictionary-hardness results in Section~\ref{sec:vac-minimization} intentionally isolate only the explicit-dictionary selection layer. Once a finite candidate pool of verifier-visible cells has already been exposed as an unrestricted explicit dictionary, selecting a smallest accepted subcover is set-cover hard; transfer of that hardness to any particular pipeline-generated family requires additional representability assumptions.

\subsection{Classical Lineage: Proof Complexity, Trusted Checking, and PL Views of Verification}

The broader theoretical lineage comes from proof complexity, model checking, and small trusted kernels. Cook and Reckhow's classical separation between semantic proof content and the overhead of a concrete proof system is the clearest ancestor of our VAC/SCC/DL split~\cite{CookReckhow1979}. Model checking provides the paradigm of local facts plus global composition under a small trusted checker~\cite{Clarke1999ModelChecking}. Proof-carrying code makes the same architectural point in a different domain: trust the small checker and the carried proof, not the whole analyzer that found it~\cite{Necula1997PCC}. The recent PL-oriented positioning of neural-network verification in \cite{Cordeiro2025PLChallenge} is especially close in spirit to our own framing.

The present paper imports these ideas into CPWL neural verification in a concrete way. Section~\ref{sec:verifier-model} fixes a deliberately small verifier interface. Section~\ref{sec:guard-exactification} shows that CPWL expression graphs admit verifier-accepted local exact cells after guard-certified exactification. Section~\ref{sec:vac-proof-languages} then separates language-relative semantic cover complexity from the size of any specific proof language, and Sections~\ref{sec:vac-bounds} and \ref{sec:vac-minimization} study quantitative bounds and minimization hardness for that semantic quantity. In that sense, the paper is neither proposing a new solver proxy nor rebranding arrangement count. It is proposing a proof-centered complexity language for what a trusted verifier must actually accept.

Taken together, these literatures explain both the motivation and the limits of our claims. We build on solver-centered verification, proof-producing infrastructures, CPWL geometry, and classical proof-complexity ideas. But the central thesis of the paper is that none of solver trace length, region count, or a particular proof syntax is the fixed-language semantic cover quantity. After fixing the verifier-facing interface, the relevant comparison object is the verifier-visible local certified cell and the accepted cover built from such cells.
